\documentclass[10pt,a4paper]{amsart}
\usepackage[margin=19mm]{geometry}
\usepackage{amsmath,amssymb,mathtools}
\usepackage[scr=rsfs,cal=euler]{mathalfa}
\usepackage{xfrac}
\usepackage[unicode,hidelinks,bookmarks=false]{hyperref}
\newcommand{\ie}{i.e.\ }
\newcommand{\cf}{cf.\ }
\newcommand{\resp}{respectively\ }
\newcommand{\Subsection}{\subsection}
\providecommand{\nobreakdash}{\nobreak\hbox{-}}
\setlength{\emergencystretch}{2em}
\multlinegap0pt
\usepackage{bm}
\usepackage{bbm}
\usepackage{dsfont}
\usepackage{xspace}
\usepackage{cancel}
\usepackage{mathtools}
\usepackage{amsthm}
\makeatletter
\@ifundefined{thm}{\newtheorem{thm}{Thm}}{}
\@ifundefined{prop}{\newtheorem{prop}[thm]{Proposition}}{}
\makeatother
\def\A{\mathbb{A}}
\def\BB{\mathcal{B}}
\def\C{\mathbb{C}}
\def\HH{\mathcal{H}}
\def\KK{\mathcal{K}}
\def\Rep{\operatorname{Rep}}
\def\id{\operatorname{id}}
\title[Tubular partitions and representations of the quantum automo]{Tubular partitions and representations of the quantum automorphism group of a homogeneous rooted tree}
\author{Nathan Brownlowe, David Robertson}
\date{Source: arXiv:2509.22201v1 (CC BY 4.0)}
\begin{document}
\maketitle

\noindent\emph{Source and licence.} Tubular partitions and representations of the quantum automorphism group of a homogeneous rooted tree. Version arXiv:2509.22201v1 (CC BY 4.0), distributed under the Creative Commons Attribution 4.0 International licence, as recorded in the arXiv licence metadata for this exact version and shown on the abstract page https://arxiv.org/abs/2509.22201v1. Full rights notice: licenses/arxiv-2509.22201-CC-BY-4.0.txt.\par\smallskip

\noindent\textit{Editorial scope.} Counted unit: the Prop whose source label is \texttt{prop:psiobj} in the article (source0.tex lines 1748--1755, source label \texttt{prop:psiobj}), delivered together with the article's own complete original proof of it (source0.tex lines 1757--1818, one \texttt{proof} environment of 62 lines). No mathematics is written, repaired or re-derived: statement and proof are the article's own text line for line. Required context retained uncounted: none. Every definition, display and label of those lines is retained unchanged. Omitted from this package and not redistributed: 1--1747, 1756--1756, 1819--2137. Nothing inside a retained excerpt is omitted or altered: every retained body is delivered exactly as the article's lines are written, so no omission receipt of any kind accompanies this package. Citations: the retained excerpts carry no citation command at all, so no bibliography entry of the article is transcribed and no reference is redistributed. Reference adaptation: none was needed and none was made; every reference inside the delivered unit resolves inside it, so no unresolved reference or citation is produced. Printed slips kept literal: no printed word, symbol, hypothesis or number of the counted statement or of its proof was altered. Layout and class: the article's own class is \texttt{amsart} with packages including amscd, amssymb, mathtools, bbold, graphics, amsmath, float, enumitem, tensor, wrapfig, tikz, color, url, hyperref; the wrapper is the lane's pinned \texttt{amsart} editorial wrapper at 10pt on A4, which restates the theorem-like environments the retained text needs and adds the article's own macro definitions that the retained text needs (\texttt{\textbackslash A}, \texttt{\textbackslash BB}, \texttt{\textbackslash C}, \texttt{\textbackslash HH}, \texttt{\textbackslash KK}, \texttt{\textbackslash Rep}, \texttt{\textbackslash id}), with extra wrapper packages (bm, bbm, dsfont, xspace, cancel, mathtools). The delivered numbering is the unit's own. Assets: no retained span contains a figure environment, a graphics-inclusion command, a PDF-page inclusion, a table or a TikZ picture; no image file is copied into this package, the package declares no asset dependency and no image file is redistributed with it. Licence: version arXiv:2509.22201v1 of this article is distributed under the Creative Commons Attribution 4.0 International licence, as recorded in the arXiv licence metadata for this exact version and shown on the abstract page https://arxiv.org/abs/2509.22201v1. A full rights notice is delivered at licenses/arxiv-2509.22201-CC-BY-4.0.txt. Characters: every retained excerpt is pure ASCII, so the delivered engine, XeLaTeX with its default Latin Modern text font, reports no missing character.
\begin{prop} \label{prop:psiobj}
	Let $u = (u_{i,j})$ be a finite-dimensional unitary representation of $\A_X$ on $\HH$. Then there is a unitary representation $\psi(u) = (\psi(u)_{(x,i),(y,j)}) \in \Rep(\A_X)$ on the Hilbert space $\C^X\otimes \HH$
	such that the matrix elements satisfy the relation
	\[
	\psi(e_x\otimes u_{i,j}) = \sum_{y\in X}\psi(u)_{(x,i),(y,j)} \otimes e_y.
	\]
	Moreover, if $v \in \Rep(\A_X)$ is a representation on $\KK$, and $T\in\Rep(\A_X)(u,v)$ is an intertwiner between $u$ and $v$, then $	\id_{\C^X}\otimes T \in \BB(\C^X\otimes\HH,\C^X\otimes\KK)$	is an intertwiner between $\psi(u)$ and $\psi(v)$.
\end{prop}

\begin{proof}
	We have
	\begin{align*}
		\sum_{y\in X} \Delta(\psi(u)_{(x,i),(y,j)}) \otimes e_y &= (\Delta\otimes\id)\psi(e_x\otimes u_{i,j}) \\
		&= (\id\otimes\psi)(\psi\otimes\id)(\id\otimes\Delta)(e_x\otimes u_{i,j}) \\
		&= \sum_{1\le k\le d_u}(\id\otimes\psi)(\psi\otimes\id)(e_x\otimes u_{i,k}\otimes u_{k,j})) \\
		&= \sum_{\substack{z\in X\\ 1\le k\le d_u}} (\id\otimes\psi)(\psi(u)_{(x,i),(z,k)}\otimes e_z \otimes u_{k,j}) \\
		&= \sum_{\substack{y,z\in X\\ 1\le k\le d_u}}\psi(u)_{(x,i),(z,k)} \otimes \psi(u)_{(z,k),(y,j)} \otimes e_y,
	\end{align*}
	and comparing tensor factors gives 
	\[
	\Delta(\psi(u)_{(x,i),(y,j)}) = \sum_{\substack{z\in X\\ 1\le k\le d_u}} \psi(u)_{(x,i),(z,k)} \otimes \psi(u)_{(z,k),(y,j)}.
	\]
	To see that $\psi(u)$ is a representation, it remains to show that $\psi(u)$ is unitary. First note that 
	\begin{align*}
		\sum_{x\in X}\psi(u^*)_{(y,i),(x,j)}\otimes e_x &= \psi(e_y\otimes (u^*)_{i,j})\\
		&=\psi(e_y\otimes u_{j,i})^*\\
		&=\sum_{x\in X}\psi(u)_{(y,j),(x,i)}^*\otimes e_x\\
		&=\sum_{x\in X}(\psi(u)^*)_{(x,i),(y,j)}\otimes e_x,
	\end{align*}
	and hence we have 
	\[
	(\psi(u)^*)_{(x,i),(y,j)}=\psi(u^*)_{(y,i),(x,j)}.
	\]
	Now it follows that for each $x,y\in X$ and $1\le i,j\le d_u$ we have
	\begin{align*}
		&\sum_{\substack{z\in X\\ 1\le k\le d_u}}\psi(u)_{(x,i),(z,k)}\psi(u)_{(z,k),(y,j)}^*\otimes e_z\\
		&\hspace{3cm}= \sum_{\substack{z\in X\\ 1\le k\le d_u}}\psi(u)_{(x,i),(z,k)}\psi(u^*)_{(y,k),(z,j)}\otimes e_z\\
		&\hspace{3cm}= \sum_{1\le k\le d_u}\Big(\sum_{z\in X}\psi(u)_{(x,i),(z,k)}\otimes e_z\Big)\Big(\sum_{z'\in X}\psi(u^*)_{(y,k),(z',j)}\otimes e_{z'}\Big)\\
		&\hspace{3cm}=\sum_{1\le k\le d_u}\psi(e_x\otimes u_{i,k})\psi(e_y\otimes u_{k,j}^*)\\
		&\hspace{3cm}=\psi\Big(e_xe_y\otimes\sum_{1\le k\le d_u}u_{i,k}u_{k,j}^*\Big)\\
		&\hspace{3cm}=\delta_{x,y}\delta_{i,j}\psi(e_x\otimes 1_{\mathbb{A}_X})\\
		&\hspace{3cm}=\delta_{x,y}\delta_{i,j} \sum_{z\in X}a_{x,z}\otimes e_z,
	\end{align*}
	and hence for each $z\in X$ we have 
	\[
	\sum_{1\le k\le d_u}\psi(u)_{(x,i),(z,k)}\psi(u)_{(z,k),(y,j)}^*=\delta_{x,y}\delta_{i,j} a_{x,z}.
	\]
	Hence 
	\[
	\sum_{\substack{z\in X\\ 1\le k\le d_u}}\psi(u)_{(x,i),(z,k)}\psi(u)_{(z,k),(y,j)}^*=\delta_{x,y}\delta_{i,j}\sum_{z\in X}a_{x,z}=\delta_{x,y}\delta_{i,j}1_{\mathbb{A}_X}
	\]
	showing that $\psi(u)$ is a unitary representation. 
	
	For the second claim, we view $T=(T_{i,j})$ and $(\id_{\C^X}\otimes T)_{(x,i),(y,j)}$ as matrices; we need to show that for any $x,y\in X$ and $1\leq i \leq d_v$, $1\leq j \leq d_u$
	\[
	\sum_{\substack{z\in X\\ 1\le k\le d_u}}(\id_{\C^X}\otimes T)_{(x,i),(z,k)}\psi(u)_{(z,k),(y,j)}=\sum_{\substack{z\in X\\ 1\le k\le d_u}}\psi(v)_{(x,i),(z,k)}(\id_{\C^X}\otimes T)_{(z,k),(y,j)},
	\]
	which is equivalent to 
	\begin{equation}\label{eq:IntIdTesorT}
		\sum_{1\le k\le d_u}T_{i,k}\psi(u)_{(x,k),(y,j)}=\sum_{1\le k\le d_u}\psi(v)_{(x,i),(y,k)}T_{k,j}.
	\end{equation}
	Equation \ref{eq:IntIdTesorT} follows from the calculation:
	\begin{align*}
		\sum_{y\in X}\Big(\sum_{1\le k\le d_u}T_{i,k}\psi(u)_{(x,k),(y,j)}\Big)\otimes e_y 
		&= \sum_{1\le k\le d_u}T_{i,k}\Big(\sum_{y\in X}\psi(u)_{(x,k),(y,j)}\otimes e_y\Big)\\
		&= \sum_{1\le k\le d_u}T_{i,k}\psi(e_x\otimes u_{k,j})\\
		&=\psi\Big(e_x\otimes \Big(\sum_{1\le k\le d_u}T_{i,k}u_{k,j}\Big)\Big)\\
		&=\psi\Big(e_x\otimes \Big(\sum_{1\le k\le d_u}v_{i,k}T_{k,j}\Big)\Big)\\
		&=\sum_{y\in X}\Big(\sum_{1\le k\le d_u}\psi(v)_{(x,i),(y,k)}T_{k,j}\Big)\otimes e_y.
	\end{align*}
\end{proof}

\end{document}
